\section{Background}
Nepal lies in an active Himalayan tectonic region, where the continuing interaction of the Indian and Eurasian plates gives rise to damaging earthquakes. Seismic demand is not spatially uniform: local geological conditions, soil deposits, topography, structural configuration, construction practice and the nature of ground shaking can influence the response experienced by individual buildings. Investigations of the Nepalese building stock show that the vulnerability of reinforced-concrete buildings is associated not only with shaking intensity but also with structural layout, detailing and local site effects \cite{gautam2016,adhikari2022}. For the assessment of an individual building, a single assumed foundation restraint may therefore conceal relevant changes in structural response.

Hospitals represent an especially important class of buildings because earthquake consequences are not limited to the possibility of structural collapse. Damage or excessive deformation can interrupt clinical services, compromise access and circulation, or affect structural and nonstructural elements needed for emergency response. Even a building that retains its overall gravity-load-carrying capacity may have insufficient functionality immediately after a major earthquake. Assessments of hospitals in Nepal and comparisons of seismic provisions for hospital structures have consequently emphasized the need to consider structural performance alongside the wider objective of post-earthquake service continuity \cite{guragain2004,subedi2022}.

Conventional analysis frequently assumes that the base of a reinforced-concrete frame is perfectly restrained. In reality, a foundation transmits forces and moments into soil that possesses finite stiffness and may deform under loading. The resulting soil--structure interaction (SSI) couples the motion of the ground, foundation and superstructure. Foundation flexibility permits additional translation and rotation, modifies the overall lateral stiffness, and can alter the natural period, participating mode shapes, lateral displacement, inter-storey drift and distribution of internal force demand \cite{wolf1985,gazetas1991}. The magnitude of the effect depends on both soil properties and the structural system; an increase in period does not necessarily imply a reduction in all seismic forces.

Several modelling approaches are available to examine SSI. Direct methods include a surrounding soil domain, whereas substructure methods replace the soil--foundation system with equivalent stiffness or impedance components acting at the base of the structural model. The Winkler idealization represents local soil resistance using springs, while elastic-half-space formulations such as those presented by Gazetas provide a basis for estimating translational, rocking and torsional foundation impedances \cite{wolf1985,gazetas1991}. These approaches make it possible to explore sensitivity to soil support while retaining a conventional three-dimensional building model. Their validity remains dependent on the assumed material properties, foundation geometry, spring formulation and modelling constraints.

For the present investigation, the G+4 Infectious and Communicable Disease Hospital in Lekhnath, Pokhara is represented by six separately identified structural blocks. The source model schedules provide fixed-base and three flexible-support conditions, classified as hard, medium and soft soil. Each block is examined across those four alternatives, allowing changes in the dynamic and response-spectrum quantities to be compared within the same structural configuration. The principal results comprise the modal periods and mass participation, lateral displacement and drift, base and storey shear, and gravity-load foundation settlement, vertical reactions and member-force envelopes.

Fragility relationships complement deterministic structural-response quantities by associating an earthquake intensity measure with a probability of exceeding a defined damage state. Their interpretation requires care because a response-spectrum maximum and a probabilistic damage-state exceedance model are not interchangeable. Likewise, adopted inter-storey drift screening intervals cannot automatically replace the thresholds used to construct an existing fragility curve \cite{giordano2021,adhikari2022,fema356}. This study therefore considers the supplied X-labelled fragility results as a separate, preliminary component of the assessment, alongside the principal comparisons of soil support and structural demand.

The significance of the investigation lies in assessing several response measures for the same hospital blocks rather than inferring performance from fundamental period alone. Differences among storey displacement, drift, foundation response and member forces are relevant to understanding how the assumption of a fixed or flexible foundation influences seismic performance assessment. These comparisons also indicate which aspects of model definition and foundation characterization warrant additional site-specific investigation before design or safety conclusions are drawn.
